\documentclass[11pt,a4paper]{article}
\usepackage[utf8]{inputenc}
\usepackage[T1]{fontenc}
\usepackage[portuguese]{babel}
\usepackage[a4paper,margin=2.5cm]{geometry}
\usepackage{booktabs}
\usepackage{amsmath}
\usepackage{verbatim}

\title{Segurança de Dados\\\large Assignment 4: Block ciphers and modes of operation}
\author{Hugo Soares}
\date{30 de outubro de 2026}

\begin{document}
\maketitle

\section*{Questão 1}

\textbf{Ideia.} Em CTR o keystream é $E_K(\text{nonce} \| \text{contador})$ e a
cifra é só $C = P \oplus K$. O AES é usado para gerar o keystream, mas a
operação final continua a ser XOR. Logo, para trocar $P$ por $P'$ sem saber
$K$, anula-se $P$ e põe-se $P'$:
$C' = C \oplus P \oplus P'$.

\textbf{Código.}

\begin{verbatim}
p  = b"PLC=07;CMD=SET;NAOH_PPM=00100;"
pp = b"PLC=07;CMD=SET;NAOH_PPM=11100;"
forged = xor(q1_ct, xor(p, pp))
\end{verbatim}

\textbf{Resultado.} Confirmei com a chave que o ciphertext dado decifra para o
valor esperado (\texttt{00100}), e a forja é:

\begin{verbatim}
Q1 forged: 51073653f3851c51af06c462d293862e93941666f8d561dbeeddc89cd972
Q1 changed bytes: [(24, 0xef, 0xee), (25, 0xdc, 0xdd)]
Q1 decrypts to: PLC=07;CMD=SET;NAOH_PPM=11100;
\end{verbatim}

Só mudam os bytes 24 e 25, que correspondem aos dois primeiros dígitos de
\texttt{00100} $\to$ \texttt{11100}: \texttt{0x30}$\oplus$\texttt{0x31} =
\texttt{0x01}, por isso \texttt{ef}$\to$\texttt{ee} e \texttt{dc}$\to$\texttt{dd}.
O contador do segundo bloco não interfere porque a alteração é feita byte a
byte dentro do mesmo bloco.

\textbf{Q1.b.} Não. A migração de ChaCha20 para AES-128-CTR não corrige a
malleability: ambos são cifras de fluxo. O AES só é usado para produzir o
keystream; a operação que combina com o plaintext continua a ser XOR e não há
MAC. Como o formato é público, o atacante sabe exatamente que bytes alterar.
O que resolve o problema é um modo autenticado (AES-GCM, ChaCha20-Poly1305) ou
encrypt-then-MAC; a escolha do algoritmo de bloco é irrelevante para esta
falha. (O nonce único continua a ser necessário, mas isso é a proteção contra
reutilização de keystream, não contra malleability.)

\section*{Questão 2}

\textbf{Ideia.} ECB cifra cada bloco de 16 bytes de forma independente e
determinística. Então dois blocos de plaintext iguais dão sempre o mesmo bloco
de ciphertext. Sem decifrar nada, basta comparar os ciphertexts bloco a bloco.

\textbf{Código.}

\begin{verbatim}
blocks = [[c[k:k+16] for k in range(0, len(c), 16)] for c in q2_cts]
for i in range(len(q2_cts)):
    for j in range(i + 1, len(q2_cts)):
        same = q2_cts[i] == q2_cts[j]
        eq = sum(a == b for a, b in zip(blocks[i], blocks[j]))
        print(f"C{i+1} vs C{j+1}: identical={same}, equal blocks={eq}")
\end{verbatim}

\textbf{Resultado.} Os quatro ciphertexts têm 64 bytes (4 blocos):

\begin{center}
\begin{tabular}{lcc}
\toprule
Par & Iguais? & Blocos iguais \\
\midrule
$C_1$ vs $C_3$ & sim & 4/4 \\
todos os outros pares & não & 0/4 \\
\bottomrule
\end{tabular}
\end{center}

Ou seja, $C_1$ e $C_3$ são exatamente o mesmo comando cifrado; $C_2$ e $C_4$
são distintos de todos os outros. Não é preciso a chave: a igualdade dos
ciphertexts prova a igualdade dos plaintexts, porque ECB é determinístico.

\textbf{O que ECB revela.} Revela igualdade de blocos, prefixos comuns e
repetições na estrutura da mensagem. Como os comandos têm formato fixo
(\texttt{PLC=<id>;CMD=SET;NAOH\_PPM=<valor>;}), dois comandos com o mesmo
prefixo de 16 bytes produzem o mesmo primeiro bloco. Confirmei isto com uma
chave de demonstração e dois valores diferentes:

\begin{verbatim}
Q2 demo C1: 3bee1322ae4c14c822683a2c86d9a400
            d5873fd250bf85cb298bebe675469a3e
Q2 demo C2: 3bee1322ae4c14c822683a2c86d9a400
            7e5a853b32bf93b30602b59db4ec980b
Q2 demo same first block: True
\end{verbatim}

Além disso, ECB permite cortar e colar blocos, reordenar blocos e montar um
dicionário de blocos frequentes. Para tráfego de PLC com mensagens muito
estruturadas isto é inaceitável: um atacante distingue comandos iguais,
deteta padrões e pode manipular blocos sem conhecer a chave.

\section*{Questão 3}

\textbf{Ideia.} Em CBC, $c_0 = E_K(P_0 \oplus iv_0)$ e, para uma mensagem
injetada, $c_1 = E_K(P \oplus iv)$. Como o IV é um contador previsível, sabemos
o IV de cada query antes de a enviar. Se fizermos $P \oplus iv = P_0 \oplus iv_0$,
então $c_1 = c_0$ e o bloco cifrado é o mesmo. Resolvendo em ordem a $P$:
$P = P_0 \oplus iv_0 \oplus iv$.

O bloco secreto tem 16 bytes e é \texttt{PLC=<id>;CMD=SET;N}; só os dois
dígitos do \texttt{<id>} são desconhecidos. São 100 hipóteses, por isso testo
cada uma com o oráculo e comparo o primeiro bloco com $c_0$. Cada query gasta
um IV novo, por isso o teste tem de usar o IV dessa query, não sempre $iv_1$.

\textbf{Código.} O oráculo é só para emular o centro de controlo:

\begin{verbatim}
def test_block(guess, iv):
    block = b"PLC=" + guess + b";CMD=SET;N"
    query = xor(block, xor(iv0, iv))   # P xor iv == block xor iv0
    return cbc_encrypt(key, iv, query)[:16] == c0

counter = int.from_bytes(iv1, "big")
matches = []
for i in range(100):
    iv = counter.to_bytes(16, "big")
    if test_block(f"{i:02d}".encode(), iv):
        matches.append(i)
    counter += 1
\end{verbatim}

\textbf{Resultado.}

\begin{verbatim}
Q3 candidate ids: ['42']
Q3 block: PLC=42;CMD=SET;N
\end{verbatim}

Só uma hipótese dá $c_0$; a verificação com a chave confirma
\texttt{PLC=42}. A lógica do ataque não usa a chave: só o oráculo a usa para
responder às queries, como o enunciado permite.

\textbf{IV aleatório.} Com IV aleatório e imprevisível, o atacante não sabe o
IV da próxima mensagem quando escolhe $P$, logo não consegue forçar
$P \oplus iv = P_0 \oplus iv_0$; o bloco de entrada de $E_K$ é
imprevisível. Um IV único mas previsível não chega: a unicidade evita
reutilização de IV entre mensagens, mas não impede este ataque de
chosen-plaintext, porque o problema é a previsibilidade. Para CBC, o IV deve
ser aleatório (ou um nonce imprevisível).

\section*{Questão 4}

\textbf{Q4.a.} O raciocínio é igual ao da Q1: a parte de confidencialidade do
GCM é CTR, por isso
$C' = C \oplus P \oplus P'$:

\begin{verbatim}
Q4 forged: 8fc0cc2e6567e77bcfaad7c8363ae39b2db91d4df9aa5acb74dde650cfe6
Q4 tag:    a903a1a5b290dfe9d41c2634deed9cb1 (inalterada)
\end{verbatim}

Mudam outra vez os bytes 24 e 25 (\texttt{0x75}$\to$\texttt{0x74},
\texttt{0xdc}$\to$\texttt{0xdd}), como na Q1. Ao decifrar e verificar, o PLC
obtém:

\begin{verbatim}
Q4 forged rejected: InvalidTag
\end{verbatim}

O GCM calcula a tag com GHASH sobre o ciphertext (e AAD/comprimentos) e
encripta o resultado com $E_K(J_0)$; a tag autentica o ciphertext inteiro.
Alterar qualquer byte muda o GHASH e a tag deixa de bater. Assim, embora a
parte CTR continue maleável, a verificação rejeita a mensagem antes de ela ser
usada. É a autenticação que faltava na Q1, não uma cifra diferente.

\textbf{Q4.b.} Reutilizar (chave, nonce) é fatal nos dois sentidos.
Na confidencialidade, o keystream é o mesmo nas duas mensagens, logo
$C_1 \oplus C_2 = P_1 \oplus P_2$; conhecendo $P_1$ (o formato é público)
recupera-se $P_2$. Em código, com uma segunda mensagem cifrada com o mesmo
nonce:

\begin{verbatim}
Q4.b recovered other: PLC=07;CMD=SET;NAOH_PPM=99999;
\end{verbatim}

Na autenticação, a tag é $T = E_K(J_0) \oplus \mathrm{GHASH}_H(\cdot)$. Com o
mesmo nonce, $E_K(J_0)$ é igual nas duas tags, por isso $T_1 \oplus T_2$ cancela
essa parcela e fica só a diferença dos GHASH. Com plaintexts conhecidos, isso
dá equações em $H$ (a chave de GHASH); recuperado $H$, o atacante calcula tags
válidas para qualquer mensagem sob aquele nonce, mesmo sem a chave AES. Portanto
o nonce tem de ser único por chave (contador ou 96 bits aleatórios); a
unicidade é o que separa o GCM de um CTR com MAC.

\section*{Conclusão}

Os modos de operação não acrescentam integridade por si: CTR continua maleável
(Q1) e CBC com IV previsível permite chosen-plaintext (Q3). ECB revela
igualdade e estrutura, logo não serve para comandos de formato fixo (Q2). A
solução é um modo autenticado (GCM) com nonce único (Q4), e nunca reutilizar
(key, nonce).

\end{document}
